=====================================================
\section{Introduction}\label{sec:intro}
% =============================================================

Automated decision-making (ADM) systems based on machine learning (ML) have become a routine
component of modern organizations. Predictive and classification models are deployed to support
tasks ranging from credit scoring and fraud detection to medical triage, customer churn, and
human-resource analytics. As these systems mature, however, an important limitation has become
increasingly evident: their numerical outputs---such as risk scores or class
probabilities---are frequently consumed in isolation, divorced from the organizational,
operational, and strategic context in which a real decision is taken
\citep{tariq2024a2c,varga2023cooperative}. A high predicted probability of an adverse event is
rarely a sufficient instruction for action; what to do with that prediction depends on broader
contextual considerations that vary across departments, time horizons, and strategic priorities.

Real-world decision-making is inherently cooperative \citep{tariq2024a2c}. On the one hand,
multiple predictive models may coexist for the same domain---for instance, departmental
classifiers tuned for sales, research and development, and human resources---and their outputs
must be combined coherently into a single, comparable risk indicator. On the other hand, experts
must overlay their judgment about which cases are most critical \emph{given the current context},
where contextual factors are often vague, gradual, and dependent on scenarios that are difficult
to encode as crisp rules. Fuzzy logic and fuzzy sets provide a mature mathematical apparatus to
represent precisely this kind of gradual, linguistically driven preference structure
\citep{kumar2025fuzzy,ghosh2020fuzzydss}.

Recently, \citet{novoa2026cademas} introduced \textbf{CADEMAS} (Cooperative Automated
Decision-Making Systems), a conceptual framework that explicitly separates three logical
concerns: an \emph{Input Manager} that ingests cases and predictive model outputs, an
\emph{Output Manager} that combines those outputs into a uniform risk score, and a
\emph{Context Model} that encodes the contextual suitability of each case via fuzzy reasoning.
A subsequent application by \citet{godz2026attrition} demonstrated the operational value of
CADEMAS in an employee-attrition prevention setting, showing how departmental classifiers can be
coordinated with contextual policies to produce case prioritizations that align with managerial
strategy.

Despite the conceptual maturity of CADEMAS, a key practical gap remains: there is no
readily available, configurable software tool that automates the framework, integrates multiple
ML predictive models, and operationalizes the fuzzy Context Model in a way that practitioners
and researchers can use without re-implementing the underlying machinery. Existing fuzzy and
MCDM packages (e.g.,\ \citealp{wieckowski2022pyfdm,wieckowski2023pyifdm,
roszkowska2024intuitionisticfuzzy,grazioso2023eligere,dikopoulou2021fcm,chen2020fuzzyr}) offer
powerful building blocks but, as discussed in Section~\ref{sec:related}, none of them is
designed as an end-to-end cooperative ADM platform that couples a weighted ensemble of ML
classifiers with a configurable fuzzy context layer.

This paper addresses that gap by presenting \textbf{CADEMAS-ML}, a software tool that
automates the CADEMAS framework for cooperative, context-aware ADM. The contributions of this
work are the following:

\begin{itemize}
\item We present CADEMAS-ML, an operational software tool that automates the CADEMAS
framework \citep{novoa2026cademas} and exposes its three logical layers (Input Manager,
Output Manager, Context Model) as configurable modules.
\item We formalize the predictive risk component as a weighted ensemble of multiple ML
classifiers, where weights are derived from user-selected performance indicators
(accuracy, F1-score, ROC-AUC) and applied to normalized model outputs to yield a global
risk score~$R_i \in [0,1]$.
\item We design the fuzzy context suitability component using configurable membership
functions (triangular, trapezoidal), categorical mappings, derived rules, and aggregation
operators (OWA, t-norm/t-conorm, weighted mean), yielding a contextual suitability
score~$S_i \in [0,1]$ that can be audited and modified by domain experts.
\item We combine the two components through a tunable trade-off
$P_i(\lambda) = \lambda, R_i + (1-\lambda), S_i$, allowing the decision-maker to
navigate the continuum between purely predictive and purely contextual rankings.
\item We illustrate the tool through an employee-attrition case study
\citep{godz2026attrition} combining departmental classifiers with two contrasting
contexts---economic crisis and digital transformation---and show that identical
predictive risks may yield different prioritization outcomes depending on context.
\end{itemize}

The remainder of the paper is organized as follows. Section~\ref{sec:background} reviews the
CADEMAS framework and related work. Section~\ref{sec:tool} describes CADEMAS-ML.
Section~\ref{sec:case} reports the employee-attrition case study.
Section~\ref{sec:conclusions} concludes.

% =============================================================
\section{Background and Related Works}\label{sec:background}
% =============================================================

\subsection{The CADEMAS Framework}\label{sec:cademas}

CADEMAS \citep{novoa2026cademas} is a framework for cooperative ADM systems. Its central
observation is that, in operational settings, a decision is rarely the output of a single
classifier applied in isolation. Instead, multiple predictive models---frequently with different
scopes, training data, and performance profiles---must be combined with contextual
considerations reflecting the strategic priorities of the organization. CADEMAS organizes
these elements into three logical layers. The \emph{Input Manager} formalizes the population
of cases under analysis and the set of available predictive models. The \emph{Output Manager}
normalizes and combines model outputs into a single risk indicator, ensuring that scores from
heterogeneous classifiers are comparable on a common scale. The \emph{Context Model}
formalizes contextual suitability through fuzzy reasoning, allowing experts to encode gradual,
scenario-dependent preferences as membership functions and aggregation rules. A final aggregator
combines the predictive and contextual components through an adjustable trade-off parameter,
yielding a global prioritization score.

The framework was recently applied in an employee-attrition prevention setting by
\citet{godz2026attrition}. In that work, departmental classifiers were combined with a fuzzy
context model encoding the cost and strategic relevance of acting on each employee,
illustrating that the same predictive risk can lead to different recommended actions under
different organizational contexts. CADEMAS bridges two largely disjoint communities---ML-based
predictive modeling and fuzzy/MCDM-based decision support---through an explicit cooperative
structure. What has been missing is an operational implementation that practitioners can
configure for new problems. CADEMAS-ML, presented in Section~\ref{sec:tool}, is designed to
fill that role.

\subsection{Literature Review on Fuzzy and Cooperative Decision Support Tools}
\label{sec:related}

We review recent work along four directions relevant to CADEMAS-ML: (a)~cooperative and
human-AI decision frameworks; (b)~fuzzy decision-support software tools; (c)~open-source fuzzy
and MCDM packages; and (d)~hybridization of ML with fuzzy or expert-knowledge components.
For each work we note whether the associated software is released as open source, since this
dimension is critical for reproducibility and practical adoption.

\paragraph{Cooperative and human-AI decision frameworks.}
The need to combine ML predictions with human and contextual judgment has motivated several
cooperative architectures. \citet{tariq2024a2c} propose an AI-to-collaborator (A2C) framework
in which humans and learning systems share decision authority through explicit interaction
protocols; the framework is described conceptually and is not released as a public software
package. \citet{varga2023cooperative} review cooperative AI architectures for high-stakes
domains and emphasize that meaningful cooperation requires shared context representations
between models and experts, but no off-the-shelf tooling is offered. CADEMAS
\citep{novoa2026cademas} extends this line by explicitly separating the predictive ensemble
from a configurable contextual policy layer, and CADEMAS-ML operationalizes that separation
as a configurable software tool.

\paragraph{Fuzzy decision-support software tools.}
A growing number of works release software artifacts to support fuzzy decision-making.
\citet{wieckowski2024makedecision} present \textsc{MakeDecision} (\textit{SoftwareX}, 2024,
doi:~10.1016/j.softx.2024.101658), a web-based tool for graphical design of decision-making
models in crisp and fuzzy environments that exposes intermediate calculations to support
auditability; it is freely accessible but its source code is not fully open.
\citet{chekry2024dematel} introduce \textsc{PyDEMATEL} (\textit{SoftwareX}, 2024,
doi:~10.1016/j.softx.2024.101889), an open-source Python package implementing DEMATEL and
fuzzy DEMATEL that emphasizes traceability of intermediate computations as a prerequisite for
trust. The \textsc{pyFDM} library \citep{wieckowski2022pyfdm} (\textit{SoftwareX}
20:101271, doi:~10.1016/j.softx.2022.101271) and its intuitionistic extension
\textsc{PyIFDM} \citep{wieckowski2023pyifdm} (\textit{SoftwareX}, 2023,
doi:~10.1016/j.softx.2023.101344) provide open-source Python implementations of fuzzy and
intuitionistic-fuzzy MCDM operators, including aggregation functions and ranking procedures.
These tools are valuable building blocks but their scope is limited to fuzzy decision-making
per se; they do not integrate ML classifiers nor formalize a cooperative ADM architecture.

\paragraph{Open-source fuzzy and MCDM packages.}
A complementary line of work focuses on releasing reusable packages.
\citet{roszkowska2024intuitionisticfuzzy} present \textsc{IFMCDM} (\textit{SoftwareX}, 2024,
doi:10.1016/j.softx.2024.101721), an Rpackage for intuitionistic fuzzy MCDM.
\citet{grazioso2023eligere} describe \textsc{Eligere} (\textit{IJIDeM}, 2023,
doi:~10.1007/s12008-023-01568-2), an open-source platform for multi-criteria decision
support in engineering design. \citet{dikopoulou2021fcm} release a Fuzzy Cognitive Maps
Python package for modeling causal contextual dependencies. All three are openly distributed
but remain methodological libraries: they do not aggregate predictive ML outputs, nor do
they expose a configurable trade-off between predictive and contextual components.

\paragraph{Hybridization of ML with fuzzy or expert knowledge.}
A third strand combines ML classifiers with fuzzy or expert components.
\citet{kumar2025fuzzy} survey hybrid fuzzy-ML architectures for uncertainty-aware DSS and
identify the need for tools that make the fuzzy layer easily configurable by domain experts.
\citet{ghosh2020fuzzydss} describe a fuzzy DSS incorporating expert rules into a predictive
workflow. \citet{suvitha2025fermateanfuzzy} extend the fuzzy toolbox with Fermatean fuzzy
operators (\textit{JIDMIS}, 2025, doi:~10.59543/jidmis.v2i.12601) suited to elicited expert
preferences. \citet{aly2018fuzzyahp} integrate fuzzy AHP with ML for context-sensitive
prioritization. These works confirm the methodological soundness of hybrid fuzzy-ML systems
but typically present bespoke implementations tied to a single domain.

Table~\ref{tab:related} summarizes representative works along these four directions.
Two observations stand out. First, a significant and growing share of recent fuzzy
decision-support tools is released as open source: all four SoftwareX entries carry
confirmed DOIs and publicly accessible code or web applications
(doi:~10.1016/j.softx.2024.101658, 10.1016/j.softx.2024.101889,
10.1016/j.softx.2022.101271, 10.1016/j.softx.2023.101344, 10.1016/j.softx.2024.101721).
Second, none of the reviewed works simultaneously provides (i)~an explicit cooperative
architecture combining multiple ML predictive models and (ii)~a configurable fuzzy context
model with an adjustable trade-off between predictive and contextual components.
CADEMAS-ML fills that intersection, building on \citet{novoa2026cademas} and
\citet{godz2026attrition}.

\begin{table}[htbp]
\centering
\small
\caption{Comparison of representative recent works on fuzzy and cooperative decision-support
tools. The Software/Openness column reflects availability as reported by the authors.
Entries marked $\dagger$ have DOIs confirmed by independent search.}
\label{tab:related}
\begin{tabularx}{\linewidth}{l X X l}
\toprule
\textbf{Work} & \textbf{Approach} & \textbf{Methodology} & \textbf{Software/Openness} \
\midrule
\citet{novoa2026cademas} & Cooperative ADM framework & Cooperative; context-aware; fuzzy & Conceptual \
\citet{godz2026attrition} & CADEMAS applied to HR attrition & Cooperative; fuzzy; ML & Application \
\citet{tariq2024a2c} & Human--AI framework (A2C) & Cooperative; hybrid & Not specified \
\citet{varga2023cooperative} & Cooperative AI review & Cooperative; survey & Not specified \
\citet{wieckowski2024makedecision}$^{\dagger}$ & \textsc{MakeDecision} web tool & Fuzzy; MCDM; auditable & Free web app \
\citet{chekry2024dematel}$^{\dagger}$ & \textsc{PyDEMATEL} & Fuzzy; DEMATEL; Python & Open-source (Python) \
\citet{wieckowski2022pyfdm}$^{\dagger}$ & \textsc{pyFDM} Python library & Fuzzy MCDM & Open-source (Python) \
\citet{wieckowski2023pyifdm}$^{\dagger}$ & \textsc{PyIFDM} Python library & Intuitionistic fuzzy MCDM & Open-source (Python) \
\citet{roszkowska2024intuitionisticfuzzy}$^{\dagger}$ & \textsc{IFMCDM} R package & Intuitionistic fuzzy MCDM & Open-source (R) \
\citet{grazioso2023eligere}$^{\dagger}$ & \textsc{Eligere} DSS platform & Fuzzy ranking; MCDM & Open-source \
\citet{dikopoulou2021fcm} & FCM Python package & FCM; context modelling & Open-source (Python) \
\citet{chen2020fuzzyr} & \textsc{FuzzyR} R package & Fuzzy inference & Open-source (R) \
\citet{kumar2025fuzzy} & Hybrid fuzzy-ML survey & Fuzzy; ML; hybrid & Survey (no code) \
\citet{ghosh2020fuzzydss} & Fuzzy DSS + expert rules & Fuzzy; hybrid & Not specified \
\citet{suvitha2025fermateanfuzzy}$^{\dagger}$ & Fermatean fuzzy operators & Fuzzy aggregation & Not specified \
\citet{aly2018fuzzyahp} & Fuzzy AHP + ML & Fuzzy MCDM; ML & Not specified \
\textbf{CADEMAS-ML} (this work) & Software tool for cooperative ADM & Coop.; fuzzy; ML; context & Open-source planned \
\bottomrule
\end{tabularx}
\end{table}

In summary, the literature confirms that fuzzy logic and cooperative architectures are
well-established responses to the limitations of purely predictive ADM. The available
software ecosystem, however, is fragmented: open-source packages provide solid fuzzy or
MCDM primitives but do not integrate ML predictive ensembles, while cooperative ADM
frameworks remain mostly conceptual. CADEMAS-ML fills this gap by providing a configurable
tool that unifies a weighted predictive ensemble with a fuzzy contextual layer.